\section{Introduzione al capitolo}\label{sec:intro_cap4}
Questo capitolo illustra nel dettaglio l'architettura tecnica e le scelte implementative che costituiscono il motore dell'LLM Query Builder.
Mentre il capitolo precedente ha descritto il sistema dal punto di vista dell'utente finale, qui verranno analizzati i componenti software, lo stack tecnologico e le metodologie di addestramento utilizzate per trasformare il linguaggio naturale in query SPARQL eseguibili.

La trattazione seguirà il flusso logico della costruzione del sistema: partendo dall'infrastruttura di base (\cref{sec:stack}), passeremo alla complessa pipeline di fine-tuning del modello (\cref{sec:finetuning}), per poi analizzare l'architettura del backend (\cref{sec:backend}) e le specifiche del dataset di addestramento (\cref{sec:dataset}).

\section{Stack tecnologico: Fedora e Blazegraph}\label{sec:stack}
L'infrastruttura dati su cui si appoggia l'LLM Query Builder è costituita dall'accoppiata Fedora Commons e Blazegraph.
Blazegraph, in particolare, opera come triplestore ad alte prestazioni per il Datavault, fornendo l'endpoint SPARQL nativo interrogato dal nostro sistema.
La sfida architettonica principale è consistita nel creare un layer intermedio capace di dialogare fluidamente con la rigidità formale richiesta da Blazegraph, astraendo completamente questa complessità dal frontend esposto al ricercatore.


\section{Pipeline di fine-tuning dell'LLM}\label{sec:finetuning}
Per ottenere traduzioni accurate dal linguaggio naturale a SPARQL, l'utilizzo di un LLM generico si è rivelato insufficiente, rendendo necessario un processo di domain adaptation tramite fine-tuning.
    \subsection{Dataset generativo e schema SPARQL}\label{subsec:dataset_gen}
    Il cuore dell'addestramento è basato su un dataset generativo, prodotto tramite lo script \code{genera\_dataset.py}.
    La versione finale (\textit{v4}) comprende 16 categorie distinte di query, progettate per coprire l'intero spettro delle interrogazioni tipiche del dominio umanistico.
    Durante la fase di validazione, l'analisi delle chiamate di rete fallite aveva evidenziato un tasso di errore significativo dovuto a una scorretta generazione della sintassi SPARQL\@.
    Nello specifico, il modello tendeva a fallire a causa di un'errata gestione dei prefissi dei metadati.
    Per risolvere la problematica in modo strutturale, si è proceduto a raffinare gli esempi all'interno del dataset di addestramento, correggendo esplicitamente il modo in cui il modello mappa e dichiara i prefissi ontologici, e gestendo puntualmente i \textit{type-mismatches} tra proprietà testuali e numeriche.

    \subsection{Configurazione dell'addestramento con Unsloth}\label{subsec:unsloth}
    Per ottimizzare i tempi di addestramento e ridurre l'impronta in memoria VRAM, la pipeline è stata implementata avvalendosi della libreria Unsloth.
    La configurazione ha previsto l'uso di precisione mista (\code{fp16} e \code{bf16}) e tecniche di \textit{gradient checkpointing}.
    Inoltre, l'impiego del \textit{data packing} ha permesso di concatenare sequenze più brevi per massimizzare l'efficienza computazionale durante l'aggiornamento dei pesi.


\section{Architettura del backend e integrazione frontend}\label{sec:backend}
    \subsection{Inference engine (FastAPI)}\label{subsec:fastapi}
    L'orchestrazione delle richieste è demandata a un backend sviluppato in Python tramite il framework FastAPI.\@
    L'inference engine è stato calibrato per garantire la consistenza del \textit{system prompt} a ogni invocazione, vincolando il modello a mantenere il ruolo di traduttore tecnico.
    Particolare attenzione è stata posta alla corretta generazione del token EOS (\textit{End Of Sentence}), fondamentale per troncare l'output del modello ed evitare allucinazioni aggiuntive in coda alla query SPARQL.\@

    Lo sviluppo iterativo di questi componenti in Python e del frontend in React è stato ampiamente facilitato dall'impiego di ambienti di sviluppo avanzati basati su agenti, come Antigravity IDE, che hanno permesso di ottimizzare la stesura del codice e testare rapidamente le configurazioni dei prompt di sistema direttamente all'interno della codebase.

    \subsection{Integrazione con Blazegraph tramite proxy Vite}\label{subsec:vite}
    Sul versante client, l'integrazione è gestita tramite un proxy configurato in Vite.
    Questa architettura previene problemi legati alle policy CORS e instrada le query SPARQL generate direttamente dall'interfaccia React all'endpoint di Blazegraph, garantendo un flusso dati bidirezionale sicuro e performante tra l'utente e il Datavault.


\section{Formato dati: dataset\_datavault\_v4.jsonl}\label{sec:dataset}
    \subsection{Le 8 regole di generazione SPARQL}\label{subsec:regole}
    Il formato dati alla base del fine-tuning, consolidato nel file \\ \code{dataset\_datavault\_v4.jsonl}, struttura la conoscenza del dominio in formato JSONL.\@
    Per istruire l'LLM, sono state codificate 8 specifiche regole di generazione SPARQL.\@
    Queste regole impongono al modello restrizioni rigorose su quali relazioni attraversare, come filtrare le date e come gestire la paginazione e i limiti dei risultati.

    \subsection{Gestione degli errori: OUT\_OF\_DOMAIN}\label{subsec:errori}
    Un requisito fondamentale per l'affidabilità in un contesto accademico è la capacità del sistema di ammettere i propri limiti.
    Nel dataset è stata implementata un'architettura di \textit{error handling} JSON-based: quando l'utente pone una domanda irrilevante rispetto allo schema del Datavault, il modello è addestrato a non tentare di forzare una query SPARQL, ma a restituire un esplicito stato \code{OUT\_OF\_DOMAIN}.
    Questo fallback previene l'esecuzione di query malformate e migliora la trasparenza dell'interfaccia utente.


\section{Conclusioni del capitolo}\label{sec:conclusioni_cap4}
Le scelte architetturali descritte in questo capitolo --- dal raffinamento del dataset per i prefissi, fino all'implementazione del backend FastAPI --- dimostrano come la transizione da un proof-of-concept teorico a uno strumento funzionante richieda un'attenta ingegnerizzazione del dato e del software.
Nel \cref{chap:valutazione} esamineremo come questa architettura si comporta concretamente in fase di test quantitativo e qualitativo.